\ssr{ЗАКЛЮЧЕНИЕ}

В ходе выполнения лабораторной работы была достигнута поставленная цель, заключающаяся в сравнении последовательной и параллельной версий алгоритма проверки связности графа после удаления заданной на вход $k$-ой вершины.

В рамках работы были выполнены следующие задачи:
\begin{enumerate}
	\item описан последовательный алгоритм проверки связности графа после удаления вершины $k$;
	\item разработана параллельная версия алгоритма проверки связности графа;
	\item реализованы последовательная и параллельная версии алгоритма на языке программирования C++;
	\item проведён сравнительный анализ зависимости времени выполнения задачи проверки связности графа после удаления вершины $k$ от размерности входных данных для последовательной версии алгоритма и для параллельной версии, запущенной с одним рабочим потоком;
	\item выполнен сравнительный анализ зависимости времени работы параллельной версии алгоритма проверки связности графа после удаления вершины $k$ от размерности входных данных при использовании различного числа вспомогательных потоков ($1, 2, 4, 8, 16, 32, 64$).
\end{enumerate}

По результатам проведённого исследования установлено, что параллельная реализация алгоритма обеспечивает преимущество по времени выполнения на графах большого размера, при этом наибольшее ускорение достигается при использовании ограниченного числа потоков, сопоставимого с количеством логических ядер вычислительной системы.
